\documentclass[12pt]{ximera}
\title{Final Exam}
\begin{document}
\begin{abstract}
The cumulative final: how modified divisors move quotas, which methods suffer which
paradoxes and quota violations, a full five-method apportionment of girl scouts, the
subtraction game, saddle points, and bonus problems on Condorcet candidates, weighted
voting, and fair division.
\end{abstract}
\maketitle

\section*{Final Exam}

\begin{problem}
Assess whether each of the following claims is true or false.

\begin{prompt}
\textbf{(a)} Claim: the quotas of states with the \emph{largest} populations will change
the fastest when the standard divisor of an apportionment problem is decreased.

\begin{multipleChoice}
\choice[correct]{True}
\choice{False}
\end{multipleChoice}

\begin{feedback}
True. A quota is $q = \frac{p}{d}$, so changing the divisor scales every quota by the
same factor --- and the same percentage change produces a larger absolute movement for a
larger quota.
\end{feedback}
\end{prompt}

\begin{prompt}
\textbf{(b)} Claim: the quotas of states with the \emph{smallest} populations will change
the fastest when the standard divisor of an apportionment problem is increased.

\begin{multipleChoice}
\choice{True}
\choice[correct]{False}
\end{multipleChoice}

\begin{hint}
Compare this carefully with part (a). Does the direction of the divisor change affect
\emph{which} states move fastest?
\end{hint}

\begin{feedback}
False. The largest states' quotas move fastest in \emph{either} direction --- increasing
the divisor does not reverse which states are most sensitive to it.

Concretely, raising the divisor from $30$ to $33$ drops a state of population $230$ from
$7.667$ to $6.970$, a change of $0.70$, while a state of population $43$ drops only from
$1.433$ to $1.303$, a change of $0.13$.
\end{feedback}
\end{prompt}

\begin{prompt}
\textbf{(c)} Claim: a \emph{zero-sum} game is a game where every point earned by one of
the players is lost by the other.

\begin{multipleChoice}
\choice[correct]{True}
\choice{False}
\end{multipleChoice}

\begin{feedback}
True --- that is the definition. The players' payoffs always sum to zero, which is
exactly why a single payoff matrix from one player's perspective captures the whole
game.
\end{feedback}
\end{prompt}

\begin{prompt}
\textbf{(d)} Claim: every 2-player zero-sum game has a \emph{saddle point}.

\begin{multipleChoice}
\choice{True}
\choice[correct]{False}
\end{multipleChoice}

\begin{feedback}
False. Problem 5(b) below is a counterexample: its maximin is $-1$ and its minimax is
$2$, so no saddle point exists. Such games have no stable pair of pure strategies, and
both players are better off mixing.
\end{feedback}
\end{prompt}
\end{problem}

\begin{problem}
Identify all applicable answers.

\begin{prompt}
\textbf{(a)} Which of the following can occur when using Hamilton's method of
apportionment?

\begin{selectAll}
\choice[correct]{Alabama paradox}
\choice{Upper quota violation}
\choice[correct]{New States paradox}
\choice{Lower quota violation}
\end{selectAll}

\begin{feedback}
The \textbf{Alabama paradox} and the \textbf{New States paradox} (and the population
paradox as well).

Hamilton's method always satisfies the quota rule --- every state receives either its
lower or its upper quota --- so neither quota violation is possible. What it cannot avoid
are the paradoxes, which arise from handing out surplus seats by comparing fractional
parts.
\end{feedback}
\end{prompt}

\begin{prompt}
\textbf{(b)} Which of the following can occur when using Webster's method of
apportionment?

\begin{selectAll}
\choice[correct]{Upper quota violation}
\choice{New States paradox}
\choice[correct]{Lower quota violation}
\choice{Population paradox}
\end{selectAll}

\begin{hint}
Webster's method is a divisor method, and it rounds conventionally rather than always up
or always down. What does each of those two facts rule in or out?
\end{hint}

\begin{feedback}
Both quota violations --- \textbf{upper} and \textbf{lower}.

Webster's method is a divisor method, so it is immune to all of the paradoxes. But
because it rounds conventionally rather than always in one direction, its bias is small
but can go either way, so it can break the quota rule at either end.

Homework 8 has an upper-quota violation: district A had a standard quota of $86.915$ ---
upper quota $87$ --- and Webster's method awarded it $88$.

Compare with the one-directional methods: Jefferson's can only violate the upper quota,
and Adams' only the lower.
\end{feedback}
\end{prompt}
\end{problem}

\begin{problem}
Twenty girl scouts are selling cookies and plan to knock on $600$ doors across four
neighborhoods, so they need to apportion scouts to each neighborhood.

\begin{center}
\begin{tabular}{|r|c|c|c|c|c|}\hline
\textbf{Neighborhood} & A & B & C & D & \textbf{Total}\\\hline
\textbf{Population} & 43 & 103 & 224 & 230 & 600\\\hline
\end{tabular}
\end{center}

\begin{prompt}
\textbf{(a)} Compute the standard divisor and standard quotas, to three decimals.

Standard divisor: $\answer{30}$

\smallskip
$q_A = \answer[tolerance=0.0005]{1.433}$ \quad $q_B = \answer[tolerance=0.0005]{3.433}$
\quad $q_C = \answer[tolerance=0.0005]{7.467}$ \quad $q_D = \answer[tolerance=0.0005]{7.667}$

\begin{feedback}
\[ SD = \frac{600}{20} = 30, \]
giving quotas $1.433$, $3.433$, $7.467$, and $7.667$, which total $20$. \checkmark
\end{feedback}
\end{prompt}

\begin{prompt}
\textbf{(b)} Complete this apportionment with Hamilton's method.

Lower quotas: $A: \answer{1}$, $B: \answer{3}$, $C: \answer{7}$, $D: \answer{7}$,
totaling $\answer{18}$, so the surplus is $\answer{2}$ scouts.

\smallskip
Hamilton's apportionment: $A: \answer{1}$ \quad $B: \answer{3}$ \quad $C: \answer{8}$
\quad $D: \answer{8}$

\begin{feedback}
The lower quotas total $1+3+7+7 = 18$, leaving $2$ surplus scouts. The fractional parts
are
\[ A: .433, \qquad B: .433, \qquad C: .467, \qquad D: .667, \]
so D and C --- the two largest --- each gain a scout:
\[ A = 1, \quad B = 3, \quad C = 8, \quad D = 8. \]
(A and B tie at $.433$, but neither is in contention for a surplus seat, so the tie never
has to be broken.)
\end{feedback}
\end{prompt}

\begin{prompt}
\textbf{(c)} Round the standard quotas as if doing just the first step of each method.

Jefferson (down): $A: \answer{1}$, $B: \answer{3}$, $C: \answer{7}$, $D: \answer{7}$,
total $\answer{18}$

\smallskip
Adams (up): $A: \answer{2}$, $B: \answer{4}$, $C: \answer{8}$, $D: \answer{8}$,
total $\answer{22}$

\smallskip
Webster (conventional): $A: \answer{1}$, $B: \answer{3}$, $C: \answer{7}$, $D: \answer{8}$,
total $\answer{19}$

\smallskip
Huntington-Hill cutoffs: $A: \answer[tolerance=0.005]{1.414}$,
$B: \answer[tolerance=0.005]{3.464}$, $C: \answer[tolerance=0.005]{7.483}$,
$D: \answer[tolerance=0.005]{7.483}$, giving $A: \answer{2}$, $B: \answer{3}$,
$C: \answer{7}$, $D: \answer{8}$ for a total of $\answer{20}$

\begin{feedback}
\[
\begin{array}{lccccl}
 & A & B & C & D & \text{total}\\
\text{quota} & 1.433 & 3.433 & 7.467 & 7.667 & 20\\
\text{Jefferson} & 1 & 3 & 7 & 7 & 18\\
\text{Adams} & 2 & 4 & 8 & 8 & 22\\
\text{Webster} & 1 & 3 & 7 & 8 & 19
\end{array}
\]

The Huntington-Hill cutoffs are
\[
\sqrt{1\cdot 2}=1.414, \quad \sqrt{3\cdot 4}=3.464, \quad \sqrt{7\cdot 8}=7.483,
\quad \sqrt{7\cdot 8}=7.483.
\]
Comparing: $1.433 > 1.414$ rounds up to $2$; $3.433 < 3.464$ rounds down to $3$;
$7.467 < 7.483$ rounds down to $7$; $7.667 > 7.483$ rounds up to $8$. That totals exactly
$20$, so Huntington-Hill needs no modified divisor.

Note that C and D share the same cutoff of $7.483$ but land on opposite sides of it ---
the method separates two neighborhoods whose quotas differ by only $0.2$.
\end{feedback}
\end{prompt}

\begin{prompt}
\textbf{(d)} Complete the apportionment using each method, choosing the appropriate
modified divisor from $28$, $29.5$, and $33$.

\textbf{Jefferson} --- divisor $\answer{28}$, giving $A: \answer{1}$, $B: \answer{3}$,
$C: \answer{8}$, $D: \answer{8}$

\smallskip
\textbf{Adams} --- divisor $\answer{33}$, giving $A: \answer{2}$, $B: \answer{4}$,
$C: \answer{7}$, $D: \answer{7}$

\smallskip
\textbf{Webster} --- divisor $\answer[tolerance=0.05]{29.5}$, giving $A: \answer{1}$,
$B: \answer{3}$, $C: \answer{8}$, $D: \answer{8}$

\begin{hint}
Part (c) tells you which way each divisor must move: Jefferson is short by $2$ and needs
a smaller divisor, Adams is over by $2$ and needs a larger one, Webster is short by $1$.
\end{hint}

\begin{feedback}
\textbf{Jefferson}, $d = 28$: modified quotas $1.536$, $3.679$, $8.000$, $8.214$;
rounding down gives $1, 3, 8, 8 = 20$. \checkmark

\textbf{Adams}, $d = 33$: modified quotas $1.303$, $3.121$, $6.788$, $6.970$; rounding up
gives $2, 4, 7, 7 = 20$. \checkmark

\textbf{Webster}, $d = 29.5$: modified quotas $1.458$, $3.492$, $7.593$, $7.797$;
conventional rounding gives $1, 3, 8, 8 = 20$. \checkmark

Comparing all five:
\[
\begin{array}{lcccc}
 & A & B & C & D\\
\text{Hamilton} & 1 & 3 & 8 & 8\\
\text{Jefferson} & 1 & 3 & 8 & 8\\
\text{Adams} & 2 & 4 & 7 & 7\\
\text{Webster} & 1 & 3 & 8 & 8\\
\text{Huntington-Hill} & 2 & 3 & 7 & 8
\end{array}
\]
The bias is on full display: Adams moves two scouts from the big neighborhoods to the
small ones, while Jefferson and Webster agree with Hamilton. Huntington-Hill sits in
between, giving the smallest neighborhood the extra scout that Adams does, but without
Adams' larger shift.
\end{feedback}
\end{prompt}
\end{problem}

\begin{problem}
Alice and Bob are playing a variation of the Subtraction Game: start with $18$ stones,
each turn a player takes $1$, $2$, or $3$ stones, and the winner takes the final stone.

\begin{prompt}
If Alice goes first, which player has the winning strategy?

\begin{multipleChoice}
\choice[correct]{Alice}
\choice{Bob}
\end{multipleChoice}

How many stones should she take first? $\answer{2}$

\begin{feedback}
Alice wins. With moves of $1$ to $3$, the losing positions are the multiples of $4$, and
$18 = 4\cdot 4 + 2$ is not one.

She takes $2$ stones, leaving $16$. From then on she answers each of Bob's moves so the
pair removes exactly $4$: if Bob takes $k$, she takes $4-k$. Bob faces $16$, $12$, $8$,
$4$, and finally $0$ --- so Alice takes the last stone.
\end{feedback}
\end{prompt}
\end{problem}

\begin{problem}
The following payoff matrices each represent Alice's perspective for a zero-sum game.
Identify the saddle point if it exists.

\begin{prompt}
\textbf{(a)}
\[
\begin{array}{c|c|c|}
\phantom{|} & c & d \\\hline
a & -7 & 2 \\\hline
b & 0 & 5 \\\hline
\end{array}
\]
Does a saddle point exist?
\begin{multipleChoice}
\choice[correct]{Yes}
\choice{No}
\end{multipleChoice}
If so, its value is $\answer{0}$.

\begin{feedback}
Row minima $-7$ and $0$ give a maximin of $0$; column maxima $0$ and $5$ give a minimax
of $0$. The saddle point is $0$ at cell $(b,c)$.
\end{feedback}
\end{prompt}

\begin{prompt}
\textbf{(b)}
\[
\begin{array}{c|c|c|}
\phantom{|} & c & d \\\hline
a & 4 & -3 \\\hline
b & -1 & 2 \\\hline
\end{array}
\]
Does a saddle point exist?
\begin{multipleChoice}
\choice{Yes}
\choice[correct]{No}
\end{multipleChoice}

\begin{feedback}
Row minima $-3$ and $-1$ give a maximin of $-1$; column maxima $4$ and $2$ give a minimax
of $2$. Since $-1 \neq 2$, there is no saddle point --- this is the counterexample for
Problem 1(d).
\end{feedback}
\end{prompt}

\begin{prompt}
\textbf{(c)}
\[
\begin{array}{c|c|c|c|}
\phantom{|} & d & e & f \\\hline
a & -1 & 5 & 0 \\\hline
b & -2 & 0 & -3 \\\hline
c & 2 & 3 & 1 \\\hline
\end{array}
\]
Does a saddle point exist?
\begin{multipleChoice}
\choice[correct]{Yes}
\choice{No}
\end{multipleChoice}
If so, its value is $\answer{1}$.

\begin{feedback}
Row minima are $-1$, $-3$, and $1$, so the maximin is $1$. Column maxima are $2$, $5$,
and $1$, so the minimax is $1$. The saddle point is $1$ at cell $(c,f)$.
\end{feedback}
\end{prompt}

\begin{prompt}
\textbf{(d)}
\[
\begin{array}{c|c|c|}
\phantom{|} & d & e \\\hline
a & 0 & -3 \\\hline
b & 2 & 4 \\\hline
c & -5 & 1 \\\hline
\end{array}
\]
Does a saddle point exist?
\begin{multipleChoice}
\choice[correct]{Yes}
\choice{No}
\end{multipleChoice}
If so, its value is $\answer{2}$.

\begin{feedback}
Row minima are $-3$, $2$, and $-5$, so the maximin is $2$. Column maxima are $2$ and $4$,
so the minimax is $2$. The saddle point is $2$ at cell $(b,d)$.

Row $b$ dominates both other rows for Alice, and column $d$ dominates $e$ for Bob, so
both players' choices are forced.
\end{feedback}
\end{prompt}
\end{problem}

\begin{problem}
Complete the payoff matrix from Alice's perspective for a two-player zero-sum game
similar to Rock/Paper/Scissors, except that only Alice has paper and only Bob has
scissors:
\begin{itemize}
\item rock wins \emph{1 point} versus scissors;
\item scissors wins \emph{2 points} versus paper;
\item paper wins \emph{1 point} versus rock;
\item everything else is a tie.
\end{itemize}

\[
\begin{array}{c|c|c|}
\phantom{|} & \text{Bob: rock} & \text{Bob: scissors} \\\hline
\text{Alice: rock} & \answer{0} & \answer{1} \\\hline
\text{Alice: paper} & \answer{1} & \answer{-2} \\\hline
\end{array}
\]

\begin{hint}
Entries are from Alice's perspective, so a win for Bob is negative.
\end{hint}

\begin{feedback}
\begin{itemize}
\item rock vs.\ rock: a tie, $0$.
\item rock vs.\ scissors: rock wins $1$ point for Alice, $+1$.
\item paper vs.\ rock: paper wins $1$ point for Alice, $+1$.
\item paper vs.\ scissors: scissors wins $2$ points for Bob, $-2$.
\end{itemize}

This game has no saddle point: the maximin is $-2$ and the minimax is $1$.
\end{feedback}
\end{problem}

\begin{problem}
The following are bonus questions from previous chapters.

\begin{prompt}
\textbf{(a)} Identify the Condorcet candidate in the following election.

\begin{center}
\begin{tabular}{c|ccc}
Number of Votes: & 8 & 6 & 4 \\\hline
1st & B & A & C \\
2nd & C & D & B \\
3rd & A & C & D \\
4th & D & B & A
\end{tabular}
\end{center}

\begin{multipleChoice}
\choice{A}
\choice{B}
\choice[correct]{C}
\choice{D}
\choice{There is no Condorcet candidate}
\end{multipleChoice}

\begin{hint}
A Condorcet candidate beats every other candidate head-to-head. With $18$ voters, a
matchup is won with at least $10$.
\end{hint}

\begin{feedback}
\textbf{C}. Checking C's head-to-head matchups:
\[
\begin{array}{lll}
\text{C vs A} & 12\text{--}6 & \text{C wins}\\
\text{C vs B} & 10\text{--}8 & \text{C wins}\\
\text{C vs D} & 12\text{--}6 & \text{C wins}
\end{array}
\]
C beats every other candidate, so C is the Condorcet candidate --- despite having only
$4$ first-place votes, the fewest of any candidate who received any. C is ranked 1st or
2nd on every single ballot.
\end{feedback}
\end{prompt}

\begin{prompt}
\textbf{(b)} Consider the weighted voting system $[10:7,5,3,1]$ with players A, B, C,
and D. How many winning coalitions does this system have?

$\answer{6}$

\begin{feedback}
A coalition wins when its weight reaches $10$. Checking every subset:
\[
\set{A,B}=12, \quad \set{A,C}=10, \quad \set{A,B,C}=15, \quad \set{A,B,D}=13,
\quad \set{A,C,D}=11, \quad \set{A,B,C,D}=16,
\]
which is $6$ winning coalitions.

Every one of them contains A --- without A the remaining players total only $9$, short of
the quota --- so A has veto power. Note also that $\set{A,D}=8$ and $\set{B,C,D}=9$ both
fall short.
\end{feedback}
\end{prompt}

\begin{prompt}
\textbf{(c)} Suppose Bandit, Chili, Bingo, and Bluey are having a picnic. Chili splits
the food into four shares that she views as equal, and the others view only the following
shares as fair:

\begin{center}
\begin{tabular}{c|ccc}
Name & Bandit & Bingo & Bluey \\\hline
Fair shares & $\set{s_1, s_2, s_3}$ & $\set{s_2, s_3}$ & $\set{s_1, s_3}$
\end{tabular}
\end{center}

Identify three different fair divisions. In every one of them, Chili receives
$s_{\answer{4}}$.

\smallskip
Division 1: Bandit $s_1$, Bingo $s_2$, Bluey $s_{\answer{3}}$

\smallskip
Division 2: Bandit $s_3$, Bingo $s_2$, Bluey $s_{\answer{1}}$

\smallskip
Division 3: Bandit $s_{\answer{2}}$, Bingo $s_3$, Bluey $s_1$

\begin{hint}
Nobody bids on $s_4$, so Chili --- the divider, for whom all four shares are equal --- must
take it. Then work out which assignments of $s_1$, $s_2$, $s_3$ satisfy the other three.
\end{hint}

\begin{feedback}
Chili is the divider and values all four shares equally, so any share is fair for her.
Since none of the other three players bid on $s_4$, Chili must take it --- otherwise
somebody would be stuck with a share they consider unfair.

That leaves $s_1$, $s_2$, $s_3$ for Bandit, Bingo, and Bluey. Bingo will take $s_2$ or
$s_3$, and Bluey will take $s_1$ or $s_3$; Bandit is flexible. Working through the cases
gives exactly three valid divisions:
\begin{enumerate}
\item Bandit $s_1$, Bingo $s_2$, Bluey $s_3$;
\item Bandit $s_3$, Bingo $s_2$, Bluey $s_1$;
\item Bandit $s_2$, Bingo $s_3$, Bluey $s_1$.
\end{enumerate}
In each case Chili takes $s_4$. Note that Bingo and Bluey can never both have $s_3$, which
is what rules out the remaining combinations.
\end{feedback}
\end{prompt}
\end{problem}

\end{document}
